\sscp{Merger of *z/*y}{06-01-02}
Another innovation shared by all Aru languages is merger of PMP
*z/*y > **y.\footnote{
Merger of *z/*y > **y also occurs in \tsc{Sula-Buru}.
This is probably a parallel development.}
This change also affects
PMP *i before vowels (/{\gap}V),
including when it occurs before a vowel due to loss of *h or *q.
\citet{Nivens1993,Nivens2017} reconstructs Proto-Aru **y as the
result of this merger.
Word initially, most languages then have fortition of **y > **ʤ,
with Dobel and Koba undergoing subsequent **ʤ > \it{s}.
East Kola retains initial **y = \it{y} in most cases,
but some lexemes have \it{y} {\tl} \it{ʤ} in variation.
Examples of PMP *z/*y/*i > **y in word-initial position are given
in \trf{06-06}.

\begin{table}\tcp{\tsc{Aru} *z/*y/*i > **y /{\#\gap}}{06-06}
%\begin{adjustbox}{width=\textwidth}
	\begin{threeparttable}\st{5pt}
	\begin{tabular}{lllllll}\lsptoprule
local gl.	&	far	&	path	&	sew	&	blow	&	count, speak	&	shark\\
PMP	&	\hr*\tbr{z}auq	&	\hr*\tbr{z}alan	&	\hr*\tbr{z}aqit	&	\hr*hə\tbr{y}up	&	\hr*\tbr{i}hap	&	\hr*q\tbr{i}hu\\
Proto-Aru	&	**\tbr{y}aw-	&	**\tbr{y}ala	&	**-\tbr{y}ay	&	**-\tbr{y}uɸ	&	**-\tbr{y}aɸ	&	**\tbr{y}u\\ \midrule
Ujir	&	{\hr\hr}\tbr{ʤ}aw	&	{\hr\hr}\tbr{ʤ}ala	&	{\hr\hr}-\tbr{ʤ}e\tbr{y}	&	{\hr\hr}-\tbr{ʤ}u	&		&	{\hr\hr}\tbr{ʤ}u\\
Kompane	&	{\hr\hr}\tbr{ʤ}aw	&	{\hr\hr}\tbr{ʤ}ala	&		&		&		&	\\
West Kola	&	{\hr\hr}\tbr{ʤ}aw	&	{\hr\hr}\tbr{ʤ}ala	&		&		&		&	\\
East Kola	&	{\hr\hr}\tbr{y}aw	&	{\hr\hr}\tbr{y}ala	&	{\hr\hr}-\tbr{y}e\tbr{y}	&	{\hr\hr}-\tbr{y}uf	&	{\hr\hr}-\tbr{y}af	&	{\hr\hr}\tbr{y}u\\
Lola	&	{\hr\hr}\tbr{ʤ}aw	&	{\hr\hr}\tbr{ʤ}ala	&		&		&		&	\\
Lorang	&	{\hr\hr}\tbr{ʤ}aw	&	{\hr\hr}\tbr{ʤ}ala	&		&		&		&	\\
Dobel	&	{\hr\hr}\tbr{s}aw	&	{\hr\hr}\tbr{s}ala	&	\hr\cg{(-sela)}	&	{\hr\hr}-\tbr{s}o	&	{\hr\hr}-\tbr{s}aw	&	{\hr\hr}\tbr{s}u\\
Koba	&	{\hr\hr}\tbr{s}aw	&	{\hr\hr}\tbr{s}ala	&		&		&		&	\\
Barakai	&	{\hr\hr}\tbr{ʤ}aw	&		&		&		&		&	\\
Karey	&	{\hr\hr}\tbr{ʤ}aw	&	{\hr\hr}\tbr{ʤ}al	&		&		&		&	\\
Batuley	&	{\hr\hr}\tbr{ʤ}a	&		&		&		&		&	{\hr\hr}\tbr{ʤ}u\\
Mariri	&	{\hr\hr}\tbr{ʤ}aw	&	{\hr\hr}\tbr{ʤ}el	&		&		&		&	\\
Manombai	&	{\hr\hr}\tbr{ʤ}aw	&	{\hr\hr}\tbr{ʤ}ala	&		&	{\hr\hr}-\tbr{ʤ}u	&		&	\\
NE. Tarangan	&	{\hr\hr}\tbr{ʤ}ɔ	&	{\hr\hr}\tbr{ʤ}ala	&		&		&		&	\\
SE. Tarangan	&	{\hr\hr}\tbr{ʤ}aw	&	{\hr\hr}\tbr{ʤ}al	&		&		&		&	\\
NW. Tarangan	&	{\hr\hr}\tbr{ʤ}ɔ	&	{\hr\hr}\tbr{ʤ}ala	&		&	{\hr\hr}-\tbr{ʤ}ô\su{†}	&	{\hr\hr}-\tbr{ʤ}ɔ\su{†}	&	{\hr\hr}\tbr{ʤ}u\\
SW. Tarangan	&	{\hr\hr}\tbr{ʤ}ôw	&	{\hr\hr}\tbr{ʤ}ala	&	{\hr\hr}-\tbr{ʤ}ɛ\tbr{y}	&	{\hr\hr}-\tbr{ʤ}ôw	&	{\hr\hr}-\tbr{ʤ}ôw	&	{\hr\hr}\tbr{ʤ}u\\
		\lspbottomrule
	\end{tabular}
		\begin{tablenotes}\tnsize
			\item[†] {Northwest Tarangan \it{-ʤô} ‘blow’ and \it{-}ʤɔ ‘count, speak’ are from Doka Timur village.}
		\end{tablenotes}
	\end{threeparttable}
%\end{adjustbox}
\end{table}

The reflexes of *z/*y > **y are more complex when word medial,
with many languages having an unconditioned split.
The most common reflexes of medial **y are \it{y,} {\0},
or loss with fronting of an adjacent vowel \tsc{(+fr)}.
Some languages have glide fortition with **y > **ʤ,
and Dobel and Koba have subsequent **ʤ > \it{s}.
The reflexes of *z/*y/*i > **y in word-medial position attested
in each language are summarised in \trf{06-07}, and examples are
given in \trf{06-08}.

\begin{table}
	\centering\tcp{Reflexes of PMP *z/*y/*i word medially in \tsc{Aru}}{06-07}
	\st{6pt}\begin{tabular}{llllll||l|lllll}\lsptoprule
Ujir	&	y	&		&		&	\tsc{+fr}	&	{\0}	&	Karey	&		&		&		&		&	{\0}	\\
Kompane	&		&		&		&	\tsc{+fr}	&	{\0}	&	Batuley	&	y	&		&		&	\tsc{+fr}	&	{\0}	\\
West Kola	&		&		&		&		&	{\0}	&	Mariri	&	y	&	ʤ	&		&	\tsc{+fr}	&		\\
East Kola	&		&		&		&	\tsc{+fr}	&	{\0}	&	Manombai	&	y	&		&		&	\tsc{+fr}	&		\\
Lola	&	y	&		&		&	\tsc{+fr}	&	{\0}	&	NE. Tarangan	&	y	&		&		&		&	{\0}	\\
Lorang	&		&	ʤ	&		&		&	{\0}	&	SE. Tarangan	&		&	ʤ	&		&		&	{\0}	\\
Dobel	&		&		&	s	&		&	{\0}	&	NW. Tarangan	&	y	&		&		&	\tsc{+fr}	&	{\0}	\\
Koba	&		&		&	s	&		&	{\0}	&	SW. Tarangan	&	y	&		&		&		&	{\0}	\\
Barakai	&		&		&		&		&	{\0}	&		&		&		&		&		&		\\
		\lspbottomrule
	\end{tabular}
\end{table}

\begin{table}\tcp{\tsc{Aru} *z/*y/*i > **y /V{\gap}V}{06-08}
\begin{adjustbox}{width=\textwidth}
	\begin{threeparttable}\st{2pt}
	\begin{tabular}{llllllll}\lsptoprule
local gl.	&	{\hr}rain	&	{\hr}ladder	&	\hp{-}{\hr}pay	&	{\hr}fresh water	&	{\hr}mango	&	{\hr}interior	&	{\hr}dugong\\
PMP	&	\hr*qu\tbr{z}an	&	\hr*haRə\tbr{z}an	&	\hp{-}\hr*ba\tbr{y}aD	&	\hr*wahiyəR	&	\hr*wa\tbr{i}	&	\hr*da\tbr{y}a	&	\hr*du\tbr{y}uŋ\\
Proto-Aru	&	**\cg{(w)}u\tbr{y}an	&	**Re\tbr{y}an	&	**-pa\tbr{y}aR	&	**wa\tbr{y}aR	&	**wa\tbr{y}[a]	&	**ra\tbr{y}a\su{†}	&	**r\cg{(iu)}\tbr{y}u\\ \midrule
Ujir	&	{\hr\hr}uan	&	{\hr\hr}rián	&	{\hr\hr}-ɸa\tbr{y}	&	{\hr\hr}wa\tbr{y}	&	{\hr\hr}wa	&	{\hr\hr}raa	&	{\hr\hr}riu, ri\tbr{y}u\\
Kompane	&	{\hr\hr}gwian	&		&		&	{\hr\hr}gwah	&		&	{\hr\hr}raa	&	\\
West Kola	&	{\hr\hr}uan	&		&		&	{\hr\hr}wah	&		&		&	\\
East Kola	&	{\hr\hr}uin	&	{\hr\hr}heen	&	{\hr\hr}-ɸeeh	&	{\hr\hr}weeh	&	{\hr\hr}wee	&	{\hr\hr}ree	&	{\hr\hr}riu, ri\tbr{y}u\\
Lola	&	{\hr\hr}wian	&		&		&	{\hr\hr}waʔ	&		&	{\hr\hr}ra\tbr{y}a	&	\\
Lorang	&	{\hr\hr}gú\tbr{ʤ}en	&		&		&	{\hr\hr}gwar	&		&	\hr\cg{(wúney)}	&	\\
Dobel	&	{\hr\hr}ʔu\tbr{s}an	&		&	{\hr\hr}-ɸa\tbr{s}ir	&	{\hr\hr}kwar	&	{\hr\hr}kwa\tbr{s}i	&	{\hr\hr}ra\tbr{s}a	&	{\hr\hr}riw\\
Koba	&	{\hr\hr}ʔu\tbr{s}an	&		&		&	{\hr\hr}ʍar	&		&	\hr\cg{(wúney)}	&	\\
Barakai	&	{\hr\hr}guon	&		&		&	{\hr\hr}gwaor	&		&		&	\\
Karey	&	{\hr\hr}guon	&		&		&	{\hr\hr}gaor	&		&		&	\\
Batuley	&	{\hr\hr}guon	&	{\hr\hr}rion	&		&	{\hr\hr}gwa\tbr{y}or	&		&	{\hr\hr}ra\tbr{y}	&	{\hr\hr}ri\\
Mariri	&	{\hr\hr}gu\tbr{ʤ}on	&		&		&	{\hr\hr}gwa\tbr{y}or	&		&	{\hr\hr}ree	&	\\
Manombai	&	{\hr\hr}gian	&	{\hr\hr}rean	&		&	{\hr\hr}g\tbr{a}yar	&		&	{\hr\hr}rea	&	\\
NE. Tarangan	&	{\hr\hr}gu\tbr{y}an	&		&		&	{\hr\hr}gwar	&		&	{\hr\hr}ra	&	\\
SE. Tarangan	&	{\hr\hr}gu\tbr{ʤ}an	&		&		&	{\hr\hr}gwar	&		&	{\hr\hr}ra	&	\\
NW. Tarangan	&	{\hr\hr}gian	&	{\hr\hr}rêan\su{‡}	&		&	{\hr\hr}gar	&	{\hr\hr}ga\tbr{y}a\su{‡}	&	{\hr\hr}ra	&	{\hr\hr}ru\\
SW. Tarangan	&	{\hr\hr}gô\tbr{y}an	&	{\hr\hr}re\tbr{y}an	&	{\hr\hr}-ɸa\tbr{y}ar	&	{\hr\hr}gar	&	{\hr\hr}ga\tbr{y}a	&	{\hr\hr}ra	&	{\hr\hr}ru\\
		\lspbottomrule
	\end{tabular}
		\begin{tablenotes}\tnsize
			\item[†]	{PMP *daya is reconstructed with the meaning ‘upriver, toward the interior’.
								Reflexes in the Aru languages mean ‘wooded area; land/interior’.
								East Kola \it{ree} has these meanings in addition to ‘weather’.}
			\item[‡]	{Northwest Tarangan \it{rêan} ‘ladder’
								and \it{gaya} ‘mango’ are from Rebi village.}
		\end{tablenotes}
	\end{threeparttable}
\end{adjustbox}
\end{table}
